\section{Discussion and Limitations}\label{sec:discussion}

\subsection{What the results say about agent MFA}

The simple version of the argument holds up. An approval is a second factor for an agent action when it lives
outside the agent's reach and names the exact effect, and the two halves do different work. Binding is the
dependable half: it removed every post-approval attack against \sys{} in our experiments, it is cheap, and AP2's
mandates and Microsoft's ADR~0030 implementation already have it. Trusted rendering is the fragile half: it makes the largest difference when the agent
lies (\RQ{2}, \RQ{4}), but only for an approver who reads it, and it surfaces details that cause harmless requests
to be refused. The attacks that succeeded against real designs were mostly rendering attacks: padding, spoofed
labels and displays that do not distinguish two effects.

The second device is the part of the MFA analogy that our measurements do not support directly. A trusted
confirmation on the agent's own host matched \sys{} in every experiment except the grant scenario, because no
experiment compromised the host's display. What the phone adds is a different trust domain for the decision and a signature that a third
party can check, and those are arguments from the threat model and the proofs rather than from attack success
rates. In our reading, the second device is worth its cost when the agent's host may be
compromised or when an approval has to be shown to someone else later, and it is not a substitute for binding or
for a card a person can read.

\subsection{Deployment guidance}

Four practical points follow. First, bind the approval to a canonical digest of every argument and execute the
stored call, not the agent's resubmission; the four runtimes in which we reproduced substitution lack exactly
this. Second, render the card from the call, type every value, escape what cannot be seen, and list every
recipient in the title; the agent's text, if shown at all, belongs in a collapsed box. Third, cap pending approvals
and lock out repeated denials, and accept that this only reduces the attacker's chances. Fourth, do not let the
agent's claim that a user asked for something decide whether to ask; in our scenarios that claim alone was enough
to skip approval for an allowed write.

\subsection{Threats to validity}

\paragraph{No human data.} Every approver in this paper is a simulated oracle or a language model. The oracles'
miss probabilities are assumptions, the exact-compare oracle is an upper bound, and the fatigue model is a
sensitivity analysis. The language-model results describe automated reviewers, which some agents already use, and
say nothing about how often people catch these attacks. A user study is the natural next step and is not attempted
here.

\paragraph{One scripted technique per scenario.} The compromised agent in the scripted suite uses one technique per
scenario, and the AgentDojo replay uses the benchmark's ground-truth calls with an honest or a forged narration. A
best-response attacker that tries every technique against each design would raise the baselines' attack success,
narrated cards in particular, and could find attacks on \sys{} that these workloads do not contain. The red-team
tests and the attack matrix cover the \sys{}-specific attacks separately.

\paragraph{Stimuli for model approvers.} The \llmAttackStimuli{} attack and \llmBenignStimuli{} benign stimuli are
small, hand-built and use placeholder bodies, which inflates false denies on the card (\cref{sec:rq4}). The models
were chosen automatically on one day, the provider can route a model identifier to different backends, and a
content filter blocked some answers. The per-model intervals in \cref{fig:llm} are wide for that reason.

\paragraph{Static and derived cells.} OpenClaw was reviewed from its released bundle, not run, and the framework
column of \cref{tab:attacks} is derived from the audit. The attack matrix was generated against an earlier
revision of \sys{} than the one used in the experiments, but after the defense changes it tests had been merged,
and both that run and the language-model run were made from working trees with uncommitted changes outside the
gateway code. The proprietary coding
agents were not audited.

\paragraph{Moving targets.} Every runtime and protocol we tested is under active development, and several cells may
change with the next release. Versions and commits are pinned in the tables and the artifact.
